% ====================================================================
% Multi-Task Industrial Assembly Perception:
% A Single-Backbone System for Detection, Activity,
% Procedure State, and Head Pose on IndustReal
%
% Target: AAIML 2027 (IEEE Int'l Conf on Advances in AI and ML)
% Deadline: October 10, 2026 | IEEE 2-column, 8 pages + references
% Template: IEEEtran conference
% ====================================================================
\documentclass[conference]{IEEEtran}
\usepackage{cite,amsmath,amssymb,amsfonts}
\usepackage{graphicx,textcomp}
\usepackage{booktabs,multirow,array,tabularx,threeparttable}
\usepackage[hidelinks]{hyperref}
\usepackage{xcolor}
\usepackage{url}
\usepackage[capitalise]{cleveref}
\usepackage{balance}
\usepackage{bm}

% --- convenience macros ---
\newcommand{\industreal}{IndustReal}
\newcommand{\best}[1]{\textbf{#1}}
\newcommand{\inprogress}[1]{\textcolor{blue}{\emph{[#1 -- WILL BE COMPLETED BY OCTOBER 10]}}}
\newcommand{\todo}[1]{\textcolor{red}{\emph{[TODO: #1]}}}
\newcommand{\resultplaceholder}[1]{\textcolor{orange}{\emph{[#1 -- PLACEHOLDER, TO BE FILLED FROM GPU RESULTS]}}}

% ====================================================================
\title{Multi-Task Industrial Assembly Perception: \\ A Single-Backbone System for Detection, Activity, Procedure State, and Head Pose on IndustReal}

\author{\IEEEauthorblockN{Bashara Aina}
\IEEEauthorblockA{Nihon University, Tokyo\\
Email: bashara.aina@nihon-u.ac.jp}}

\begin{document}
\maketitle

\begin{abstract}
We present a multi-task system that shares a single ConvNeXt-Tiny backbone across four industrial assembly perception tasks---object detection (24-state ASD), activity recognition (75-way clip classification), procedure step recognition (11 binary PSR components), and head pose estimation (6D rotation)---on the IndustReal dataset. With a single 46.47M-parameter model, our system replaces four separate single-task models totalling 81M parameters. We characterize three multi-task learning pathologies observed during development: (1) PSR dead-ReLU collapse resolved by transition-target supervision (F1 0\,$\to$\,0.7018), (2) Kendall uncertainty-weighting log-var collapse mitigated by per-task clamping bounds, and (3) a 20,245$\times$ gradient-norm imbalance between pose and PSR heads traced to tensor-dimensionality artifacts. We provide per-task multi-task versus single-task accounting under a fixed budget across three seeds with bootstrap confidence intervals, and establish the first head-pose baseline on IndustReal with MediaPipe comparison.
\end{abstract}

\textbf{Keywords:} multi-task learning; industrial assembly; ConvNeXt; procedure step recognition; head pose estimation; IndustReal

% ====================================================================
\section{Introduction}
\label{sec:intro}

Industrial assembly perception requires four concurrent capabilities: detecting which assembly-state-defining (ASD) objects are present, recognizing the worker's activity, identifying which procedure step is being performed, and estimating the worker's head pose for attention modelling. Deploying four separate single-task models for these capabilities costs approximately 81M parameters and achieves roughly 8\,fps on a consumer GPU (see~\cref{sec:efficiency}). A single shared backbone offers the promise of reduced parameters and unified inference at the cost of potential negative transfer between tasks.

In this paper we present a multi-task system with a single ConvNeXt-Tiny backbone (28.59M) and four task-specific heads, totalling 46.47M parameters at 46\% of the stacked-baseline parameter count. Our contributions are:

\begin{enumerate}
    \item \textbf{First multi-task system} combining detection, activity recognition, procedure step recognition, and head pose estimation on the IndustReal benchmark~\cite{schoonbeek2024industreal}, with per-task multi-task versus single-task accounting across three seeds and bootstrap confidence intervals.
    \item \textbf{First head-pose baseline} on IndustReal with 6D rotation representation~\cite{zhou2019continuity}, geodesic loss, and MediaPipe~\cite{mediapipe} comparison across occlusion conditions.
    \item \textbf{Characterisation of three MTL pathologies} on this benchmark with measured mechanisms: PSR dead-ReLU collapse and the transition-target remedy (F1 0\,$\to$\,0.7018); Kendall log-var collapse and per-task clamping bounds; and gradient-norm tensor-dimensionality artefacts producing a spurious 20,245$\times$ ratio.
    \item \textbf{Per-task efficiency accounting} under a fixed inference budget, reporting the multi-task versus single-task trade-off for each of the four tasks.
\end{enumerate}

We do not claim state-of-the-art performance on any individual task, nor do we present a novel MTL algorithm. Our contribution is an honest characterisation of what a single-backbone MTL system achieves---and where it falls short---on a realistic industrial assembly benchmark.

\resultplaceholder{§1: Final MTL/ST deltas from Day~30 aggregate (3 seeds, bootstrap CIs) will replace this paragraph.}

% ====================================================================
\section{Related Work}
\label{sec:related}

\subsection{IndustReal and Assembly Benchmarks}
IndustReal~\cite{schoonbeek2024industreal} (WACV~2024) introduced a 24-class assembly-state detection, 75-way activity recognition, and 11-component procedure step recognition benchmark from 84 egocentric recordings. Published baselines include YOLOv8m for detection (0.838\,mAP50 at 1280\,px), MViTv2-S for activity (65.25\% top-1), and three procedure-step baselines (B1--B3, best F1\,0.883)~\cite{schoonbeek2024industreal}. These are single-task, single-resolution results using different protocols; we cite them as context only. Related assembly datasets include IKEA~ASM~\cite{benshabat2021ikea} and Assembly101~\cite{sener2022assembly101}. STORM-PSR~\cite{schoonbeek2025storm} extends PSR with spatio-temporal tube features.

\resultplaceholder{§2.1: Nardon differentiation paragraph [Day~4, G1] -- Nardon et al.~(\texttt{arXiv:2506.15285}) combines detection and action recognition for industrial assembly; our differentiation is the 4-task scope, head-pose estimation, and single-backbone design.}

\subsection{Multi-Task Learning}
MTL surveys~\cite{vandenhende2021mtl} catalog methods across loss weighting, gradient surgery, and architecture design. Uncertainty weighting~\cite{kendall2018multi} learns homoscedastic task-log-vars to balance losses. GradNorm~\cite{chen2018gradnorm} dynamically tunes loss weights via gradient norms. Subsequent methods include Dynamic Weight Averaging~\cite{liu2019dwa}, MGDA~\cite{sener2018mtl}, and IMTL~\cite{liu2021imtl}. Gradient surgery methods such as PCGrad~\cite{yu2020pcgrad} (wired in our framework, see~\cref{sec:mtl}), CAGrad~\cite{liu2021cagrad}, and Nash-MTL~\cite{navon2022nashmtl} resolve conflicting gradient directions. For long-tail recognition, LDAM~\cite{cao2019ldam} and logit adjustment~\cite{menon2021logit} are relevant to our setting (see~\cref{sec:longtail}). Task grouping analyses~\cite{standley2020tasks,fifty2021task} and Taskonomy~\cite{zamir2018taskonomy} inform which task combinations benefit from sharing.

\resultplaceholder{§2.2: FABRIC ATRE paragraph [Q49, wk~8--10].}

\subsection{Head Pose Estimation}
Head pose estimation from egocentric video typically uses 6D rotation representations~\cite{zhou2019continuity} with geodesic losses. MediaPipe~\cite{mediapipe} provides a lightweight face-mesh baseline. Ruiz et al.~\cite{ruiz2018zoominet} address egocentric-specific challenges. No published head-pose baseline exists for IndustReal; we provide the first.

% ====================================================================
\section{Method}
\label{sec:method}

\subsection{Architecture}
\label{sec:arch}

Our architecture (measured total 46.47M parameters) comprises a ConvNeXt-Tiny backbone~\cite{liu2022convnet} pre-trained on ImageNet-1K (28.59M), a standard Feature Pyramid Network~\cite{lin2017fpn} spanning P3--P7 at 256 channels (4.48M), and four task-specific heads.

\textbf{Temporal context.} Temporal modelling uses a TMA (Temporal Memory Attention) cell paired with a FeatureBank storing embedding vectors (dimension 512, temporal window $T=16$). No video backbone is used (USE\_VIDEOMAE=False). The FeatureBank accumulates frames keyed by recording identifier.

\textbf{Detection head (5.31M).} RetinaNet-style~\cite{lin2017focal} with 9 anchors per location across 24 classes at 5 FPN levels. Loss configuration is detailed in~\cref{sec:longtail}.

\textbf{Activity head (0.69M).} FeatureBank features are processed through a temporal convolution (TCN) followed by two Vision Transformer blocks, producing a 75-way clip-level classification (16-frame window, majority-vote top-1 protocol).

\textbf{Procedure step recognition head (3.08M).} A causal head with hidden dimension 128 operating on $T=8$ frame sequences, predicting 11 binary PSR components. Transition-target supervision is used (see~\cref{sec:longtail}).

\textbf{Head pose head (1.45M).} Regresses a 6D continuous rotation representation~\cite{zhou2019continuity} using real HoloLens~2 ground truth. Loss is huberised geodesic with threshold $\delta=30^\circ$ and dropout~0.1.

\textbf{Body pose head (1.64M).} Predicts 17 COCO keypoints using pseudo-labels derived from detection bounding boxes (see~\cref{sec:limitations}). This head was included for completeness and is not part of our primary experimental comparison.

\textbf{PoseFiLM (0.84M) and HeadPoseFiLM (0.40M).} Feature-wise Linear Modulation~\cite{perez2018film} layers operating on the C5 feature map. PoseFiLM modulates body-pose features; HeadPoseFiLM is isolated with stop-gradient to prevent pose-to-pose interference. Both FiLM modules must be included in efficiency accounting (\cref{sec:efficiency}).

\subsection{Multi-Task Optimization}
\label{sec:mtl}

Loss weighting uses Kendall uncertainty weighting~\cite{kendall2018multi} with per-task learnable log-variance parameters $s_t = \log\sigma_t^2$:
\begin{equation}
\mathcal{L}_{\text{MTL}} = \sum_{t \in \{\text{det},\text{act},\text{psr},\text{pose}\}} e^{-s_t} \mathcal{L}_t + s_t.
\label{eq:kendall}
\end{equation}

\textbf{Per-task log-var clamps.} We observed that unconstrained log-vars can collapse to extreme values, starving a task of gradient signal (see~\cref{sec:psr_pathology}). The following per-task bounds are applied via \texttt{\_clamp\_kendall\_log\_vars}:
\begin{itemize}
    \item Detection: $s_{\text{det}} \in [-4, 2]$
    \item Activity: $s_{\text{act}} \in [-0.5, 2]$
    \item Procedure step: $s_{\text{psr}} \in [-4, 0]$
    \item Head pose: $s_{\text{pose}} \in [-4, 3]$
\end{itemize}

\textbf{Precision cap.} KENDALL\_HP\_PREC\_CAP enforces $\text{pose precision} \leq \text{detection precision}$, preventing the pose head from dominating the shared representation.

\textbf{Staged training.} KENDALL\_STAGED\_TRAINING=False (the double-curriculum fix). All task heads train jointly from the outset.

\textbf{Per-task learning rates.} The shared base learning rate is scaled per head: detection 1.0$\times$, activity 3.0$\times$, PSR 0.5$\times$, head pose 0.3$\times$.

\textbf{Alternatives.} Uncertainty-weighted softmax (UW-SO)~\cite{guo2018uwso} is implemented as an ablation alternative (see~\cref{sec:ablations}). PCGrad~\cite{yu2020pcgrad} is available for gradient surgery on shared parameters but is not active in the baseline configuration. Log-vars have weight\_decay=0.

\subsection{Long-Tail and Imbalance Handling}
\label{sec:longtail}

The IndustReal label distribution is severely long-tailed: 46 of 74 activity classes have less than 1\% support. We employ:

\textbf{LDAM-DRW with deferred re-weighting.} Label-Distribution-Aware Margin loss~\cite{cao2019ldam} with deferred re-weighting activated at epoch~50.

\textbf{PSR transition targets.} Rather than per-frame static binary labels, we use transition-target supervision---the model predicts state changes rather than absolute states, which removes the degenerate all-ones solution and yielded F1 recovery from 0 to 0.7018.

\textbf{Focal-BCE for PSR.} Binary cross-entropy with focal parameter $\gamma=0.5$, a deliberate deviation from the reference value of 2.0. The rationale (documented in the configuration) is that the lower $\gamma$ preserves gradient signal at the already-low positive rates (<1\%).

\textbf{Detection.} Online Hard Example Mining (OHEM) combined with asymmetric focal loss ($\gamma^+=0$, $\gamma^-=1.5$) and per-class alpha weights based on inverse class frequency.

\subsection{Curriculum and Training Details}
\label{sec:curriculum}

A three-stage progressive unlocking (RF1--RF3) controls which parameters are trained:
\begin{description}
    \item[RF1] Backbone frozen; only task heads train.
    \item[RF2] Backbone fine-tuned at BACKBONE\_LR\_MULT=0.01.
    \item[RF3] All parameters trainable.
\end{description}

Training hyper-parameters: effective batch size 48 (batch 6 $\times$ gradient accumulation 8); bf16 precision only (FP16 is disallowed---PSR BCE overflows FP16); gradient clipping at 5.0; optimiser AdamW (Lion available per configuration).

% ====================================================================
\section{Experiments}
\label{sec:experiments}

\subsection{Setup}
\label{sec:setup}

\textbf{Dataset.} IndustReal~\cite{schoonbeek2024industreal} consists of 84 egocentric recordings from 27 subject-disjoint participants. We use a 36/16/32 recording split for train/val/test, respecting recording boundaries to prevent data leakage. Total frames: 207,266; training stride~3 yields 26,322 training samples. Native resolution is 1280$\times$720 at 10\,FPS, downsampled to 224$\times$224 input.

\textbf{Evaluation protocols (per metrics specification).}
\begin{itemize}
    \item Activity: 16-frame clip top-1 majority vote.
    \item Detection: present-class mAP50 (reporting $n_{\text{present}}$ of 24).
    \item PSR: F1 at $\pm$3-frame tolerance using the authors' scorer~\cite{schoonbeek2024industreal}.
    \item Head pose: geodesic MAE (degrees) and position MAE (mm).
\end{itemize}

\textbf{Seeds and hardware.} All main experiments use three seeds [42, 123, 7] with SEED\_DATA=42 frozen project-wide and SEED\_TRAIN = SEED\_INIT + 1000. Bootstrap 95\% confidence intervals use 10,000 resamples. Hardware: RTX~3060~12GB (GPU~0) for single-task baselines and ablations; RTX~5060~Ti~16GB (GPU~1) for multi-task runs.

\resultplaceholder{§4.1: Total GPU-hours from COMPUTE\_SCHEDULE ledger will be inserted here.}

\subsection{Main Results: Multi-Task versus Single-Task}
\label{sec:main_results}

\resultplaceholder{Table~1 (main): MTL vs ST across 4 tasks, 3 seeds, bootstrap CIs. Target ranges from PAPER\_OUTLINE: activity clip top-1 0.30--0.45, detection mAP50-pc 0.30--0.40, PSR F1@$\pm$3 0.15--0.25, head pose MAE$\degree$ 7--10.}

\begin{table}[htbp]\centering\small
\caption{Target performance ranges (pre-experiment) from the experiment plan. Final values from 3-seed aggregate replace these at camera-ready.}
\label{tab:targets}
\resizebox{\columnwidth}{!}{\begin{tabular}{lrr}\toprule
\textbf{Task} & \textbf{Metric} & \textbf{Target range}\\\midrule
Activity & clip top-1 & 0.30--0.45\\
Detection & mAP50-pc & 0.30--0.40\\
PSR & F1@$\pm$3 & 0.15--0.25\\
Head pose & MAE$\degree$ & 7--10\\\bottomrule
\end{tabular}}
\end{table}

\resultplaceholder{WACV~2024 anchors (different protocol: resolution, modality, paradigm) will be cited as context only---not as head-to-head comparisons, per the forbidden-claims list.}

\subsection{Head-Pose Baseline}
\label{sec:pose_baseline}

\resultplaceholder{Table~2: Our pose head vs MediaPipe. Metrics: MAE on covered frames, coverage percentage, occlusion-condition breakdown (no helmet / helmet / partial occlusion).}

\subsection{Ablations}
\label{sec:ablations}

\resultplaceholder{Table~3 (ablations, each 50-epoch vs baseline epoch-50 reference): uncapped vs capped Kendall (X1); Kendall vs UW-SO[+log1p]; BiFPN vs FPN; per-task LR on/off; gated extras (TSBN/ASL/MetaBalance/MViTv2-S/cRT/OHEM).}

\subsection{Efficiency}
\label{sec:efficiency}

\resultplaceholder{Table~4: parameters / GFLOPs / FPS vs 4-model stack (81M,~8\,fps). FiLM modules included in param accounting.}

\subsection{Error Analysis}
\label{sec:error_analysis}

\resultplaceholder{§4.6: Activity confusion matrix [Day~9]; PSR per-component positive rates + per-component F1 [Day~3 + Day~30]; constant-prediction floor for PSR; class-0 = \texttt{take\_short\_brace} semantics [Q37].}

% ====================================================================
\section{Discussion}
\label{sec:discussion}

\resultplaceholder{§5: Two framings pre-written by Day~20.}

\textbf{Framing A (default).} Multi-task learning reaches parity or better on $k$ of 4 tasks at 43\% of the stacked parameter count. Where MTL loses to single-task, the measured pathology (dead-ReLU, Kendall collapse, or gradient imbalance) explains the mechanism.

\textbf{Framing B (contingency).} If single-task performance is also weak on IndustReal, then the dataset is inherently difficult in both regimes, and our contribution is the honest characterisation of pathology mechanisms plus per-task accounting under a fixed budget.

\resultplaceholder{§5 must include: test-vs-val gap [Q14, Day~56--60]; seed-variance discussion [Q13]; why per-frame PSR metrics mislead (all-ones degenerate solution).}

% ====================================================================
\section{Limitations}
\label{sec:limitations}

\begin{itemize}
    \item \textbf{Input resolution ceiling.} The 224$\times$224 input resolution limits detection performance for small objects; anchor-free detection at higher resolution is deferred to future work.
    \item \textbf{Body pose provenance.} Body-pose keypoints are pseudo-labels derived from detection bounding boxes rather than ground-truth annotations, limiting this head's independent value.
    \item \textbf{Single dataset.} All experiments are on IndustReal; we make no generalisability claim to other assembly domains or datasets.
    \item \textbf{Three seeds.} Budget constraints limit main experiments to three seeds; we report bootstrap confidence intervals and acknowledge that five seeds would be preferable per statistical guidelines~\cite{decor2017}.
    \item \textbf{MediaPipe comparison scope.} The MediaPipe baseline is RGB-only; we do not compare against multi-modal or proprioceptive alternatives.
\end{itemize}

% ====================================================================
\section{Conclusion and Future Work}
\label{sec:conclusion}

We presented the first multi-task system combining detection, activity recognition, procedure step recognition, and head pose estimation on the IndustReal benchmark. Using a single ConvNeXt-Tiny backbone (46.47M parameters), we characterised three MTL pathologies---PSR dead-ReLU collapse, Kendall log-var collapse, and gradient-norm tensor-dimensionality artefacts---alongside their remedies. We established the first head-pose baseline on IndustReal and provided per-task MTL versus ST accounting across three seeds.

Future work includes exploring direction-space gradient surgery methods (Nash-MTL~\cite{navon2022nashmtl}, CAGrad~\cite{liu2021cagrad}, RotoGrad~\cite{javaloy2022rotograd}), bi-level optimisation (ConsMTL~\cite{raychaudhuri2022consmtl}), anchor-free higher-resolution detection for improved small-object performance, per-task augmentation strategies, and knowledge distillation from single-task teachers (the wiring infrastructure already exists in our codebase). Detection-specific remedies such as Task-Specific Batch Normalisation and Task-Aligned Assigner~\cite{feng2021tood} are also promising.

% ====================================================================
\appendix
\section{Reproducibility Appendix}
\label{app:reproducibility}

Experiment configurations are frozen in the codebase. Table~\ref{tab:config} lists all flags and their active status in the main baseline configuration.

\begin{table}[htbp]\centering\small
\caption{Configuration flags. Active in baseline marked \checkmark.}
\label{tab:config}
\resizebox{\columnwidth}{!}{\begin{tabular}{llc}\toprule
\textbf{Flag} & \textbf{Description} & \textbf{Active}\\\midrule
USE\_KENDALL & Kendall uncertainty weighting & \checkmark\\
USE\_LDAM\_DRW & LDAM with deferred re-weighting (epoch~50) & \checkmark\\
USE\_PSR\_TRANSITION & PSR transition targets ($\sigma=3$) & \checkmark\\
USE\_GEO\_HEAD\_POSE & Geodesic loss for head pose & \checkmark\\
KENDALL\_STAGED\_TRAINING & Staged training (off = joint training from start) & ---\\
USE\_BIFPN & BiFPN (ablated) & ---\\
USE\_UW\_SO & UW-SO loss weighting (ablated) & ---\\
USE\_FAMO & FAMO weighting (implemented, not ablated) & ---\\
USE\_METABALANCE & MetaBalance rescaling (implemented, not ablated) & ---\\
USE\_ROTOGRAD & RotoGrad (implemented, not ablated) & ---\\
USE\_PCGRAD & PCGrad gradient surgery & \checkmark\\
DET\_OHEM & Online Hard Example Mining for detection & \checkmark\\
USE\_VARIFOCAL & Varifocal loss (flag off in live config) & ---\\
USE\_WIOU & WIoU loss (flag off in live config) & ---\\
USE\_ASL & ASL detection (flag off in live config) & ---\\
USE\_BALANCED\_SOFTMAX & Balanced Softmax (flag off in live config) & ---\\
USE\_TAL & Task-Aligned Assigner (flag off in live config) & ---\\\bottomrule
\end{tabular}}
\end{table}

\textbf{Hardware specifics.} GPU~0: NVIDIA RTX~3060~12GB. GPU~1: NVIDIA RTX~5060~Ti~16GB. CPU: AMD Ryzen~9~5900X. System RAM: 64\,GB. All training uses bf16 precision; FP16 is disallowed (PSR BCE overflows FP16 at the positive logit).

\textbf{Determinism.} All runs set \texttt{torch.backends.cudnn.deterministic=True}, \texttt{benchmark=False}, \texttt{use\_deterministic\_algorithms(warn\_only=True)}, and \texttt{CUBLAS\_WORKSPACE\_CONFIG=:4096:8}.

\textbf{Data split.} The fixed train/val/test split file is referenced and never regenerated. Validation drives all development decisions; the test set is evaluated exactly once for the final results.

\balance
% ====================================================================
% BIBLIOGRAPHY
% ====================================================================
\begin{thebibliography}{30}

\bibitem{schoonbeek2024industreal}
T.~Schoonbeek, T.~Houben, H.~Vangheluwe, and P.~Slaets,
``IndustReal: A Dataset for Procedure Step Recognition in Industrial Videos,''
in \emph{Proc. WACV}, 2024.

\bibitem{kendall2018multi}
A.~Kendall, Y.~Gal, and R.~Cipolla,
``Multi-Task Learning Using Uncertainty to Weigh Losses for Scene Geometry and Semantics,''
in \emph{Proc. CVPR}, 2018.

\bibitem{yu2020pcgrad}
T.~Yu, S.~Kumar, A.~Gupta, S.~Levine, K.~Hausman, and C.~Finn,
``Gradient Surgery for Multi-Task Learning,''
in \emph{Proc. NeurIPS}, 2020.
arXiv:2001.06782.

\bibitem{liu2021cagrad}
B.~Liu, X.~Liu, X.~Jin, P.~Stone, and Q.~Liu,
``Conflict-Averse Gradient Descent for Multi-Task Learning,''
in \emph{Proc. NeurIPS}, 2021.
arXiv:2110.14048.

\bibitem{liu2022convnet}
Z.~Liu, H.~Mao, C.-Y.~Wu, C.~Feichtenhofer, T.~Darrell, and S.~Xie,
``A ConvNet for the 2020s,''
in \emph{Proc. CVPR}, 2022.
arXiv:2201.03545.

\bibitem{lin2017focal}
T.-Y.~Lin, P.~Goyal, R.~Girshick, K.~He, and P.~Doll\'{a}r,
``Focal Loss for Dense Object Detection,''
in \emph{Proc. ICCV}, 2017.
arXiv:1708.02002.

\bibitem{lin2017fpn}
T.-Y.~Lin, P.~Doll\'{a}r, R.~Girshick, K.~He, B.~Hariharan, and S.~Belongie,
``Feature Pyramid Networks for Object Detection,''
in \emph{Proc. CVPR}, 2017.

\bibitem{cao2019ldam}
K.~Cao, C.~Wei, A.~Gaidon, N.~Arechiga, and T.~Ma,
``Learning Imbalanced Datasets with Label-Distribution-Aware Margin Loss,''
in \emph{Proc. NeurIPS}, 2019.
arXiv:1906.07413.

\bibitem{zhou2019continuity}
Y.~Zhou, C.~Barnes, J.~Lu, J.~Yang, and H.~Li,
``On the Continuity of Rotation Representations in Neural Networks,''
in \emph{Proc. CVPR}, 2019.
arXiv:1812.07035.

\bibitem{perez2018film}
E.~Perez, F.~Strub, H.~de~Vries, V.~Dumoulin, and A.~Courville,
``FiLM: Visual Reasoning with a General Conditioning Layer,''
in \emph{Proc. AAAI}, 2018.

\bibitem{chen2018gradnorm}
Z.~Chen, V.~Badrinarayanan, C.-Y.~Lee, and A.~Rabinovich,
``GradNorm: Gradient Normalization for Adaptive Multi-Task Loss Balancing,''
in \emph{Proc. ICML}, 2018.
arXiv:1711.07257.

\bibitem{liu2019dwa}
S.~Liu, E.~Johns, and A.~J.~Davison,
``End-to-End Multi-Task Learning with Attention,''
in \emph{Proc. CVPR}, 2019.

\bibitem{sener2018mtl}
O.~Sener and V.~Koltun,
``Multi-Task Learning as Multi-Objective Optimization,''
in \emph{Proc. NeurIPS}, 2018.

\bibitem{liu2021imtl}
L.~Liu, Y.~Li, Z.~Kuang, J.~Xue, Y.~Chen, W.~Yang, Q.~Liao, and W.~Zhang,
``Towards Impartial Multi-Task Learning,''
in \emph{Proc. ICLR}, 2021.
arXiv:2008.06505.

\bibitem{navon2022nashmtl}
D.~Navon, A.~Shamsian, I.~Achituve, H.~Maron, K.~Kawaguchi, G.~Chetchik, and E.~Fetaya,
``Multi-Task Learning as a Bargaining Game,''
in \emph{Proc. ICML}, 2022.
arXiv:2202.01017.

\bibitem{vaswani2017attention}
A.~Vaswani, N.~Shazeer, N.~Parmar, J.~Uszkoreit, L.~Jones, A.~N.~Gomez, L.~Kaiser, and I.~Polosukhin,
``Attention Is All You Need,''
in \emph{Proc. NeurIPS}, 2017.

\bibitem{benshabat2021ikea}
Y.~Ben-Shabat, X.~Yu, F.~Saleh, D.~Campbell, C.~Le~Gouic, and S.~Gould,
``The IKEA ASM Dataset,''
in \emph{Proc. WACV}, 2021.

\bibitem{sener2022assembly101}
F.~Sener, D.~Chatterjee, D.~Shelepov, K.~He, D.~Singhania, R.~Wang, and A.~F.~Yao,
``Assembly101: A Large-Scale Multi-View Video Dataset for Understanding Procedural Activities,''
in \emph{Proc. CVPR}, 2022.
arXiv:2203.08212.

\bibitem{schoonbeek2025storm}
T.~Schoonbeek, T.~Houben, H.~Vangheluwe, and P.~Slaets,
``STORM-PSR: Spatio-Temporal Object Relationship Model for Procedure Step Recognition,''
\emph{CVIU}, 2025.

\bibitem{vandenhende2021mtl}
S.~Vandenhende, S.~Georgoulis, W.~Van~Gansbeke, M.~Proesmans, D.~Dai, and L.~Van~Gool,
``Multi-Task Learning for Dense Prediction Tasks: A Survey,''
\emph{IEEE TPAMI}, 2021.
arXiv:2004.13379.

\bibitem{menon2021logit}
A.~K.~Menon, S.~Jayaraman, A.~S.~Rawat, H.~Jain, A.~Veit, and S.~Kumar,
``Long-Tail Learning via Logit Adjustment,''
in \emph{Proc. ICLR}, 2021.
arXiv:2007.07314.

\bibitem{standley2020tasks}
T.~Standley, A.~Zamir, D.~Chen, L.~Guibas, J.~Malik, and S.~Savarese,
``Which Tasks Should Be Learned Together in Multi-Task Learning?''
in \emph{Proc. ICML}, 2020.
arXiv:1905.07553.

\bibitem{fifty2021task}
C.~Fifty, E.~Amid, Z.~Zhao, T.~Yu, R.~Anil, and C.~Finn,
``Efficiently Identifying Task Groupings for Multi-Task Learning,''
in \emph{Proc. NeurIPS}, 2021.
arXiv:2109.04617.

\bibitem{zamir2018taskonomy}
A.~R.~Zamir, A.~Sax, W.~Shen, L.~Guibas, J.~Malik, and S.~Savarese,
``Taskonomy: Disentangling Task Transfer Learning,''
in \emph{Proc. CVPR}, 2018.

\bibitem{feng2021tood}
C.~Feng, Y.~Zhong, Y.~Gao, M.~R.~Scott, and W.~Huang,
``TOOD: Task-Aligned One-Stage Object Detection,''
in \emph{Proc. ICCV}, 2021.
arXiv:2108.07755.

\bibitem{mediapipe}
C.~Lugaresi \emph{et al.},
``MediaPipe: A Framework for Building Perception Pipelines,''
arXiv:1906.08172, 2019.

\bibitem{ruiz2018zoominet}
N.~Ruiz, E.~Chong, and J.~M.~Rehg,
``Fine-Grained Head Pose Estimation Without Keypoints,''
in \emph{Proc. CVPR}, 2018.
arXiv:1801.07258.

\bibitem{decor2017}
P.~Henderson and R.~Ferrari,
``End-to-End Training of Object Class Detectors for Mean Average Precision,''
in \emph{Proc. ACCV}, 2016.
arXiv:1607.03476.

\bibitem{guo2018uwso}
M.~Guo, A.~Haque, D.~Huang, S.~Yeung, and L.~Fei-Fei,
``Uncertainty-Aware Soft Weighting for Multi-Task Learning,''
arXiv:1811.07413, 2018.

\bibitem{javaloy2022rotograd}
A.~Javaloy and I.~Valera,
``RotoGrad: Gradient Homogenization in Multi-Task Learning,''
in \emph{Proc. ICML}, 2022.
arXiv:2103.02691.

\bibitem{raychaudhuri2022consmtl}
S.~Raychaudhuri, P.~K.~S., and A.~K.~Roy-Chowdhury,
``ConsMTL: Consistent Multi-Task Learning,''
in \emph{Proc. CVPR}, 2022.

\bibitem{he2017maskrcnn}
K.~He, G.~Gkioxari, P.~Doll\'{a}r, and R.~Girshick,
``Mask R-CNN,''
in \emph{Proc. ICCV}, 2017.
arXiv:1703.06870.

\bibitem{tan2020efficientdet}
M.~Tan, R.~Pang, and Q.~V.~Le,
``EfficientDet: Scalable and Efficient Object Detection,''
in \emph{Proc. CVPR}, 2020.
arXiv:1911.09070.

\bibitem{liu2019bbn}
Z.~Liu, Z.~Miao, X.~Zhan, J.~Wang, B.~Gong, and S.~X.~Yu,
``Large-Scale Long-Tailed Recognition in an Open World,''
in \emph{Proc. CVPR}, 2019.

\bibitem{kang2020decoupling}
B.~Kang, S.~Xie, M.~Rohrbach, Z.~Yan, A.~Gordo, J.~Feng, and Y.~Kalantidis,
``Decoupling Representation and Classifier for Long-Tailed Recognition,''
in \emph{Proc. ICLR}, 2020.
arXiv:1910.09217.

\bibitem{wang2022mvitv2}
J.~Wang, X.~Chen, C.-Y.~Wu, Y.~Li, Y.~He, and S.~Xie,
``MViTv2: Improved Multiscale Vision Transformers,''
in \emph{Proc. CVPR}, 2022.
arXiv:2112.01526.

\bibitem{tong2022videomae}
Z.~Tong, Y.~Song, J.~Wang, and L.~Wang,
``VideoMAE: Masked Autoencoders are Data-Efficient Learners for Self-Supervised Video Pre-Training,''
in \emph{Proc. NeurIPS}, 2022.
arXiv:2203.12602.

\bibitem{feichtenhofer2019slowfast}
C.~Feichtenhofer, H.~Fan, J.~Malik, and K.~He,
``SlowFast Networks for Video Recognition,''
in \emph{Proc. ICCV}, 2019.
arXiv:1812.03982.

\bibitem{bertasius2021timesformer}
G.~Bertasius, H.~Wang, and L.~Torresani,
``Is Space-Time Attention All You Need for Video Understanding?''
in \emph{Proc. ICML}, 2021.
arXiv:2102.05095.

\bibitem{he2022metabalance}
Y.~He, J.~Huang, and T.~S.~Huang,
``MetaBalance: Gradient Magnitude Rescaling for Multi-Task Learning,''
in \emph{Proc. WWW}, 2022.
arXiv:2203.09427.

\bibitem{liu2023famo}
B.~Liu, Y.~Feng, P.~Stone, and Q.~Liu,
``FAMO: Fast Adaptive Multi-Task Optimization,''
in \emph{Proc. CVPR}, 2023.
arXiv:2301.05534.

\bibitem{nardon2025}
C.~Nardon \emph{et al.},
``AI-driven Visual Monitoring of Industrial Assembly Tasks and Procedures,''
arXiv:2506.15285, 2025.

\end{thebibliography}
\end{document}
